%-------------------------------------------------------------------------------
%	SECTION TITLE
%-------------------------------------------------------------------------------
\cvsection{Experience}


%-------------------------------------------------------------------------------
%	CONTENT
%-------------------------------------------------------------------------------
\begin{cventries}
%---------------------------------------------------------
  \cventry
    { Radboud-Universitätsklinikum Nimwegen, Humangenetik } % Organization
    { Identification and ranking of p63 binding sites putatively involved in the etiology of non-syndromic cleft lip with or without cleft palate } % Job
    { Nimwegen, Niederlande } % Location
    { November 2015 - 2016 } % Date
    { 
      \begin{cvitems} % Description(s) of tasks/responsibilities
\item { In-silico-Vorhersage und in-vitro-Validierung klinisch relevanter Transkriptionsfaktor-Bindungsstellen mittels eines integrativen Multiomics-Ansatzes. }
\item { Übernahm die Leitung eines Nassalborforschungsprojekts und verwandelte es in ein bioinformatisches Forschungsprojekt, wobei ich mir die erforderlichen bioinformatischen Fähigkeiten selbst beibrachte. }
\item { Entwicklung einer Pipeline zur Integration öffentlich verfügbarer und eigener Multi-Omics-Daten (Chip-seq, SNP, GWAS, Konservierung, Kopplungsungleichgewicht), um reproduzierbare Ergebnisse zu erzielen. haplotplotR: http: //tinyurl.com/4cjw7s5h }
\item { Effektive Kommunikation und Präsentation bioinformatischer Methoden und Ergebnisse an Forschende ohne bioinformatische Fachkenntnisse. }
      \end{cvitems}
    }

%---------------------------------------------------------
%---------------------------------------------------------
  \cventry
    { Zentrum für Molekulare und Biomolekulare Informatik CMBI, Vergleichende Genomanalyse } % Organization
    { CTCF-motif directionality controls CTCF-mediated chromatin interactions and correlates with topological domain structure } % Job
    { Nijmegen, The Netherlands } % Location
    { 2016 } % Date
    { 
      \begin{cvitems} % Description(s) of tasks/responsibilities
\item { Leitete und verwandelte ein exploratives Forschungsprojekt in ein hypothesengesteuertes Projekt, das in einer Publikation resultierte (siehe Veröffentlichungen). }
\item { Hypothesengenerierung durch Nutzung von Multi-Omics-Datensätzen, die verschiedene Dimensionen des Genoms beschreiben, einschließlich 1D (Sequenz), 2D (ChIP-Seq) und 3D-Daten (ChIA-PET, HI-C). }
\item { Hypothesentestung durch Anwendung parametrischer und nicht-parametrischer Methoden, Randomisierung sowie Modellierung von Chromatinschleifen. }
\item { Anwendung unüberwachter maschineller Lerntechniken, z.B. PCA, HCA, sowie eine Vielzahl von Visualisierungen zur Datenerkundung. }
      \end{cvitems}
    }

%---------------------------------------------------------
\end{cventries}
\newpage